% Reviewed FF.33--FF.35, integrated in F_section.tex after two independent audits. Do not input this duplicate.
\begin{claim}{FF.33}{Proved}\label{clm:FFcyclotomiccohomology}
Let \(K/U_p\) be finite. Then the integral and
rational \(p\)-Kummer kernels of \(K^\times\)
are both \(\mu_\infty\). If \(M/\mathbb Q_p\)
is finite and
\(H=\ker(\chi_p|_{P_M})\), restriction
\[
 H^1(G_M,\mathbb Q_p(1))\longrightarrow
                  H^1(H,\mathbb Q_p(1))
                                                        \tag{FF.33.1}
\]
is injective. The same is true with inertia
\(I_M\) in place of \(G_M\).
\end{claim}
\begin{proof}
Write \(K=L_0U_p\) with \(L_0/K_p\) finite,
and enlarge \(L_0\) to contain a specified
\(a\in K^\times\). Both \(T_p/K_p\) and
\(K_p(\mu_{p^\infty})/K_p\) are abelian, so
\(K/L_0\) is Galois abelian. Its cyclotomic
character has infinite image.
If \(a\in\bigcap_n(K^\times)^{p^n}\), choose
a compatible sequence of roots in \(K\).
Such a sequence exists because \(K\) contains
all roots of unity and each finite root fibre
is nonempty; the transition maps can equally
be made surjective by adjusting by roots of
unity. Its integral Kummer cocycle factors
through \(\Gamma=\operatorname{Gal}(K/L_0)\).
For a \(\mathbb Q_p(1)\)-valued cocycle \(c\)
of an abelian group, commutativity gives
\[
 (\chi_p(s)-1)c(t)=(\chi_p(t)-1)c(s).
                                                        \tag{FF.33.2}
\]
Choose \(s\) with \(\chi_p(s)\ne1\). Equation
(FF.33.2) makes \(c\) a coboundary over
\(\mathbb Q_p(1)\). Thus \(\kappa_p(a)\)
is rationally zero already over \(L_0\), so
some \(p\)-power of \(a\) belongs to the
integral kernel \(\mu_{p'}\) of
Claim~\ref{clm:FFstrictdetection}. Hence
\(a\in\mu_\infty\). The converse is immediate,
and the rational kernel is the same because
a power of an element that is a root of unity
can only come from a root of unity.

For (FF.33.1), \(P_M/H\) is central in
\(G_M/H\): the cyclotomic character embeds
it in \(\mathbb Z_p^\times\), and conjugation
preserves that character. Choose a central
\(s\in P_M/H\) with \(\chi_p(s)\ne1\).
The same cocycle computation (FF.33.2), now
using only centrality of \(s\), shows that
\(H^1(G_M/H,\mathbb Q_p(1))=0\).
Inflation--restriction gives injectivity,
because \(H\) acts trivially on the coefficient.
The identical argument applies to \(I_M/H\).
No triviality of the tame action is assumed.
\end{proof}


\begin{claim}{FF.34}{Proved}\label{clm:FFcyclotomicsplitting}
Let \(A\) be a semiabelian variety over a finite extension
\(M/\mathbb Q_p\), put
\(H=\ker(\chi_p|_{P_M})\), and set
\[
 S=\operatorname{Hom}_H(V_pA,\mathbb Q_p(1)).
                                                        \tag{FF.34.1}
\]
After a finite extension of the finite local model, the
inertia action on this space has a canonical decomposition
\[
 S=S_0\oplus S_1,\qquad
 I_M|_{S_0}=1,\quad I_M|_{S_1}=\chi_p.
                                                        \tag{FF.34.2}
\]
Both summands are stable under the full local Galois group;
its unramified action on either summand need not be trivial.
For an algebraic model over a finite extension \(K/U_p\),
the summand \(S_1\) vanishes on the rational Kummer classes
of all points over all finite extensions of \(K\).
The summand \(S_0\) is jointly detected by finitely many
such point evaluations. If \(j:X\to A\) is pointed and
its geometric image generates \(A\), those evaluations
may all be taken at points \(j(x)\).
\end{claim}
\begin{proof}
We use the convention in which \(\mathbb Q_p(1)\)
has Hodge--Tate weight \(1\). The Hodge--Tate comparison
for a smooth variety with a smooth normal-crossing
compactification applies to a semiabelian variety.
Precisely, \citet[Proposition in Section~3.4 and
Theorem~3.6, especially proof~(d)]{FF_Beilinson2012}
gives the filtered, Galois-compatible de Rham
comparison for this nonproper variety; taking
the associated graded gives the Hodge--Tate
decomposition.
Here \(H^1_{\mathrm{\acute et}}(A_{\overline M},\mathbb Q_p)
=\operatorname{Hom}(V_pA,\mathbb Q_p)\); dualizing
its logarithmic Hodge filtration gives the stated
sign convention for the Tate module.
It makes \(V_pA\) Hodge--Tate with weights \(0,1\).
One can also read the weights from the torus and abelian
quotient after the comparison theorem supplies the
decomposition: the torus has weight \(1\), and the
abelian quotient has weights \(0,1\).
Consequently \(\operatorname{Hom}(V_pA,\mathbb Q_p(1))\)
and its \(G_M\)-stable subrepresentation \(S\) are
Hodge--Tate with weights \(0,1\).
Here closure under subobjects, duals, and tensor
products, with exactness on the weight spaces, is
\citet[Theorem~4.2]{FF_Fontaine2004}.

The quotient \(P_M/H\) is central in \(G_M/H\), as in
Claim~\ref{clm:FFcyclotomiccohomology}. The tame quotient
\(I_M/P_M\) is procyclic of pro-prime-to-\(p\) order.
A central extension with procyclic quotient is abelian,
so \(I_M/H\) is abelian. Its pro-prime-to-\(p\)
subgroup has finite image in \(\operatorname{GL}(S)\):
a compact linear group has an open pro-\(p\) subgroup,
whose intersection with a pro-prime-to-\(p\) image is
trivial. Kill this finite image and the finite torsion
in the cyclotomic character by a finite extension.
On a sufficiently small open subgroup of inertia,
the action is therefore
\[
           \rho_S(g)=\exp\bigl(\log\chi_p(g)\,B\bigr)
              \quad\text{for some }B\in\operatorname{End}_{\mathbb Q_p}(S).
                                                        \tag{FF.34.3}
\]
The same formula on \(S\oplus\mathbb Q_p(1)\)
shows that the inertia Lie algebra is the
one-dimensional space generated by \(\operatorname{diag}(B,1)\).
By \citet[Theorem~3.1, with \(d=0\)]{FF_Ohkubo2014},
the Sen operator belongs to the scalar extension of
this inertia Lie algebra to \(\mathbb C_p\).
Its component on \(\mathbb Q_p(1)\) is \(1\), so
it is exactly \(\operatorname{diag}(B,1)\).
The Hodge--Tate decomposition makes that operator
semisimple, with eigenvalues \(0,1\) on \(S\).
This is the precise criterion in
\citet[Section~2.7, printed p.~45]{FF_Fontaine2004},
with the same sign convention.
Thus \(B\) is semisimple over \(\mathbb Q_p\),
with those eigenvalues. Its two eigenspaces give
(FF.34.2), after shrinking inertia once more if necessary.
Conjugation by \(G_M\) fixes the normalized generator
\(B\), since it fixes the cyclotomic character on
the central group \(P_M/H\). Hence the two eigenspaces
are \(G_M\)-stable. This argument makes no assertion
that the unramified factors have finite image.

All finite changes of local model used here can be
realized by finite extensions of the chosen algebraic
model, by the henselization/completion equivalence.
For fixed \(S\), the finite inertia images can even
be killed inside the fixed field of \(H\), since
\(H\) acts trivially on \(S\oplus\mathbb Q_p(1)\).
Enlarging the finite model inside \(K\) in this way
does not change \(K\). Alternatively one may make
a finite extension of \(K\), as allowed throughout
the point-function argument.

We prove the asserted invisibility before detection.
For \(\lambda\in S_1\), the formula for the natural
action on Hom,
\[
        g\lambda=\chi_p(g)\lambda\rho_A(g)^{-1},
\]
and (FF.34.2) say that, after choosing a basis of
the target, the underlying linear map
\(\lambda':V_pA\to\mathbb Q_p\) is inertia-equivariant
for the trivial target action. After splitting the
torus by a further finite extension, it kills its
Tate module: that module is a sum of copies of
\(\mathbb Q_p(1)\), and inertia has infinite
cyclotomic image. Thus \(\lambda'\) factors through
the Tate module of the abelian quotient \(B_A\).

Every point in question is algebraic over \(\mathbb Q\).
It is therefore defined over a finite extension \(L/K_p\),
after adjoining its finitely many coordinates to a
finite field of definition. Make a further finite
extension so that \(B_A\) has semistable reduction.
The finite-part construction used in
Claim~\ref{clm:FFpdetector} supplies a connected
\(p\)-divisible group \(\mathcal G^0\) over the
complete strict valuation ring, whose formal group
is the formal completion of the identity of the
connected N\'eron model. Restrict \(\lambda'\) to
\(V_p\mathcal G^0\) and multiply by a power of \(p\)
to make it integral. The resulting map
\(T_p\mathcal G^0\to\mathbb Z_p\) is equivariant
over the completion of \(L\). By
\citet[Theorem~4 and Corollary~1]{FF_Tate1967}, it
extends uniquely to a map of \(p\)-divisible groups
\[
                  \mathcal G^0\longrightarrow\mathbb Q_p/\mathbb Z_p.
\]
This map is zero: every finite connected source
has zero map to a constant \'etale finite group.
Thus \(\lambda'\) kills \(V_p\mathcal G^0\).

For any \(a\in B_A(L)\), some nonzero integer
multiple of \(a\) lies in that formal completion.
Indeed, the component group is finite, and every
point of the connected reduction over
\(\overline{\mathbb F}_p\) has finite order:
it and the variety descend to a finite field.
Choose compatible successive \(p\)-division points
inside the formal group. Their differences under
Galois lie in \(T_p\mathcal G^0\), so applying
\(\lambda'\) kills the full compatible Kummer
cocycle, not just its classes at separate finite
levels. Division by the integer multiple is
legitimate with rational coefficients. Restriction
to the group of any finite extension of \(K\)
therefore gives zero. This proves invisibility
of \(S_1\).

For \(\lambda\in S_0\setminus\{0\}\), inertia
equivariance permits application of
Claim~\ref{clm:FFpdetector} over a finite extension
of \(K_p\). It gives a point over a further
finite extension whose image under \(\lambda\)
has nonzero rational Kummer class. The restriction
from this inertia group to its
\(\ker(\chi_p|_P)\) subgroup is injective by
Claim~\ref{clm:FFcyclotomiccohomology}, so the
class remains nonzero after adjoining \(U_p\).
The pointed-generation assertion follows from
the corresponding assertion in
Claim~\ref{clm:FFpdetector}. Finally, a finite
family of evaluations suffices on the
finite-dimensional space \(S_0\): successively
choose an evaluation nonzero on the remaining
common kernel until that kernel is zero.
\end{proof}

\begin{claim}{FF.35}{Proved}\label{clm:FFcyclotomicreconstruct}
For every prime \(p\), every finite extension of
\[
 U_p=\operatorname{Frac}(\mathbb Z_{(p)}^{sh})
       \bigl(p^{1/m}:(m,p)=1,\ \mu_{p^\infty}\bigr)
                                                        \tag{FF.35.1}
\]
satisfies \'etale reconstruction. Every subfield of
such an extension also satisfies \'etale reconstruction.
In particular, the failure of the functional detector
in Claim~\ref{clm:FFcyclotomicdetector} does not prevent
reconstruction of \(U_5\).
\end{claim}
\begin{proof}
We prove the actual-point-function condition of
\citet[Theorem~4.18(5)]{FF_CHW2024}. Let \(K/U_p\)
be finite, let \(X/K\) be regular, connected and
quasiprojective with \(x_0\in X(K)\), and let
\(\varphi:X_{\mathrm{\acute et}}\to
\mathbb G_{m,K,\mathrm{\acute et}}\) be pinned.
Use the notation \(b(x),b_0,\xi,j:X\to A\) of
Claim~\ref{clm:FFstrictreconstruct}; in particular
\(j(x_0)=0\). As there, \(X\) is geometrically
integral and its rational points are Zariski dense.
The density proof uses only Hensel's lemma and
infinite valuation balls, so applies to the
nondiscrete henselian field \(K\).
The zero-dimensional case is already covered by
Claim~\ref{clm:FF1}.

We may pass to finite extensions of \(K\) during
this proof: a unit realizing the original geometric
point function over such an extension descends by
Galois invariance and equality of functions on a
reduced scheme, exactly as at the beginning of
Claim~\ref{clm:FFstrictreconstruct}. This does not
presuppose reconstruction over the larger field.
All finite fields considered below remain finite
extensions of \(U_p\).

\emph{1. The part detected by actual point values.}
Choose a number field \(D_0\subset K\) containing
all the finite model data and \(b_0\), with
\(K=D_0U_p\), and set \(M_0=(D_0)_v\).
Then
\[
 G_K=H_0=\ker(\chi_p|_{P_{M_0}})
                  \triangleleft G_{M_0}.
                                                        \tag{FF.35.2}
\]
The geometric part of \(\xi-\kappa_p(b_0)\)
is an integral functional \(\lambda\) in the
space \(S\) of Claim~\ref{clm:FFcyclotomicsplitting}.
After the finite changes allowed there, write
\(\lambda=\lambda_0+\lambda_1\), with
\(\lambda_i\in S_i\). For every finite-extension
point, the actual-value identity becomes
\[
  (\lambda_0)_*\kappa_A(j(x))
                =\kappa_p(b(x)/b_0)
               \quad\text{over its field of definition}.
                                                        \tag{FF.35.3}
\]
Indeed, \(\lambda_1\) kills that entire rational
Kummer class by Claim~\ref{clm:FFcyclotomicsplitting}.

Choose finitely many point evaluations jointly
detecting \(S_0\), over one finite extension \(E/K\).
Include the points, their actual values, and a
generator of \(E/U_p\) in a number field \(D_1\).
Then \(E=D_1U_p\), \(M_1=(D_1)_v\), and
\(H_1=\ker(\chi_p|_{P_{M_1}})\) is normal in
\(G_{M_1}\). Joint detection survives further
finite restriction, by restriction--corestriction
with rational coefficients.
For \(g\in G_{M_1}\), the difference
\(g\lambda_0-\lambda_0\) lies in \(S_0\):
both \(S\) and its weight-zero summand are
\(G_{M_0}\)-stable. The restricted point and
value cocycles in (FF.35.3) are invariant under
the conjugation action of \(G_{M_1}\), since
their points and actual values are defined over
\(M_1\). The same cocycle computation as
(FF.26.4) therefore makes every chosen evaluation
of \(g\lambda_0-\lambda_0\) zero. Joint detection
gives
\[
                \lambda_0\in
             \operatorname{Hom}_{G_{M_1}}(V_pA,\mathbb Q_p(1)).
                                                        \tag{FF.35.4}
\]
No equivariance of \(\lambda_1\) under the full
local Galois group is needed.

There is a normalized rational global class
\(\xi_0\) over \(X_{M_1}\) with geometric part
\(\lambda_0\), and with value zero at \(x_0\).
Here is an explicit construction, so no integrality
of the projection is silently assumed. The point
\(x_0\) splits the arithmetic fundamental-group
sequence. On
\(\pi_1(X_{\overline M_1})\rtimes G_{M_1}\),
define the cocycle by applying \(\lambda_0\)
to the geometric abelian pro-\(p\) component
and assigning zero to the section \(G_{M_1}\).
Equivariance (FF.35.4) is exactly the cocycle
identity. Equivalently, pull the universal rational
Kummer class on \(A\) back by \(j\), and apply
\(\lambda_0\). Its evaluations satisfy (FF.35.3).
The same construction defines the compatible
class over the algebraic field \(E\).
We henceforth replace \(K\) by \(E\), as permitted.

\emph{2. Boundary residues are rational.}
The summand \(\lambda_1\) is zero on the torus
of the generalized Albanese, as proved in
Claim~\ref{clm:FFcyclotomicsplitting}. Thus the
boundary residues of \(\xi_0\) are the residues
of the original integral class \(\xi\), and
belong to \(\mathbb Z_p\). We now prove that
they also belong to \(\mathbb Q\).

On a punctured rational line over a finite model,
write the normalized class in the form (FF.26.1),
with residues \(a_i\in\mathbb Z_p\). Choose
a uniformizer \(\pi\) of the finite local model
in its dense number subfield. Evaluate at
\(r_i+\pi^N\) and \(r_i+\pi^{N+1}\), for \(N\)
larger than the valuations of all nonzero root
differences. Include these points and their
actual values in a further finite number field,
with completion \(M_2\). Over
\(H_2=\ker(\chi_p|_{P_{M_2}})\), equality
(FF.35.3) identifies each evaluation with the
Kummer class of its actual value. Both classes
come from \(G_{M_2}\); the restriction injection
of Claim~\ref{clm:FFcyclotomiccohomology} therefore
gives that equality already over \(G_{M_2}\).
Apply the ordinary normalized local valuation
to the difference of the two equalities. The
other puncture factors have constant valuation,
so
\[
        a_i v_p(\pi)
           =v_p\bigl(b(r_i+\pi^{N+1})\bigr)
                -v_p\bigl(b(r_i+\pi^N)\bigr)
                    \in\mathbb Q.
                                                        \tag{FF.35.5}
\]
Consequently \(a_i\in\mathbb Q\cap\mathbb Z_p
=\mathbb Z_{(p)}\).
This calculation uses the finite local field,
not a valuation map on Kummer cohomology over
\(K\); the latter would lose this information
because \(v(K^\times)=\mathbb Q\).

For a smooth curve, use the finite \'etale map
to a punctured line with a chosen boundary
point as its only pole, exactly as in Step~1
of Claim~\ref{clm:FFstrictreconstruct}.
Corestriction has actual point values given by
products of field norms, and the residue at
infinity is the chosen residue. The rational
class is defined over a finite local field,
and its norm values are still algebraic over
\(\mathbb Q\), so the preceding comparison
applies to that class without change.
Transverse curves through the smooth locus
of each boundary component, with intersection
multiplicity one, give the higher-dimensional
assertion. Therefore the full residue divisor
\(D\) of \(\xi_0\) has coefficients in
\(\mathbb Z_{(p)}\).

\emph{3. An algebraic unit up to powers and roots of unity.}
Apply the Picard 1-motive calculation (FF.26.5)
to the rational invariant geometric class
\(\lambda_0\). Choose a positive integer \(q\)
such that \(q\lambda_0\) is integral and
\(qD\in B_X\). Such an integer exists:
first clear the finitely many denominators
of \(D\); then use
\[
 (B_X\otimes\mathbb Q_p)\cap
       \Bigl(\bigoplus_i\mathbb QD_i\Bigr)
                         = B_X\otimes\mathbb Q
\]
and clear its rational lattice denominators;
finally multiply by a power of \(p\) to make
the entire geometric class integral.
An invariant integral lift with lattice
coordinate exactly \(qD\) gives
\([qD]\in\bigcap_n p^nP(M_1)\), by the
finite-level division-point argument in
Step~3 of Claim~\ref{clm:FFstrictreconstruct}.
That intersection is finite prime-to-\(p\)
torsion. Hence a positive multiple \(dD\)
is principal. The principal divisor descends
to the number-field model and has a function
\(h\in\mathcal O(X)^\times\), normalized by
\(h(x_0)=b_0^d\).

The zero-residue difference between
\(d\lambda_0\) and the geometric Kummer class
of \(h\) lies in \((V_pP)^{G_{M_1}}=0\).
Normalization at \(x_0\) therefore gives the
rational equality
\[
                   d\xi_0=\kappa_p(h/b_0^d).
                                                        \tag{FF.35.6}
\]
Evaluate and use (FF.35.3). The rational
Kummer kernel in every finite extension of
\(K\) is \(\mu_\infty\), by
Claim~\ref{clm:FFcyclotomiccohomology}. Thus
\[
                b(x)^d=\zeta_xh(x),
                         \qquad \zeta_x\in\mu_\infty\subset K.
                                                        \tag{FF.35.7}
\]
We have proved this pointwise relation; we
do not identify the original class \(\xi\)
with the Kummer class of \(h\). Its invisible
summand may be nonzero, as Claim~\ref{clm:FFcyclotomicdetector}
demonstrates.

\emph{4. Closed points of arbitrarily large degree.}
Every finite extension \(L/K\) has infinite
absolute Galois group. Indeed,
Claim~\ref{clm:FFcyclotomictest} identifies it
with a nonabelian subgroup of the free wild
inertia group. An infinite profinite group
has open subgroups of arbitrarily large
index; otherwise its finite quotients would
have bounded order, forcing the group to
be finite. Finite separable extensions of
\(L\) therefore have unbounded degrees.
By the primitive-element theorem their
generators give closed points on
\(\mathbb A^1_L\) of unbounded degrees.
A finite nonconstant map from the proper
normalization of a geometrically integral
curve to \(\mathbb P^1_L\) then gives
closed points of unbounded degrees in
every nonempty open of that curve: choose
base points of larger degree than every
point in the finite image of its omitted
locus. Points above them have at least
the same degree. This replaces the
Eisenstein-polynomial construction, which
would not be valid for the divisible
value group of \(K\).

\emph{5. Recovering the actual point function.}
The remaining algebraic steps of
Claim~\ref{clm:FFstrictreconstruct}, Steps~4--7,
now apply, with \(\lambda_0\) in place of
\(\lambda\) and \(\mu_\infty\) in place of
\(\mu_{p'}\). We spell out why their exact
inputs have all been retained.

If \(\lambda_0\ne0\), the unit \(h\) in
(FF.35.6) is nonconstant. Relation (FF.35.7)
and the existence of large-degree points
imply that the actual point function is
nonconstant. On the nonempty inverse-image
open of \(T-1\ne0\), apply Steps~1--3 to
that target function. This gives
\((b(x)-1)^e=\eta_xg(x)\), with
\(\eta_x\in\mu_\infty\subset K\).
Use a smooth geometrically integral curve
on which restriction is injective on the
finite span of \(h^ig^j\), as supplied by
\citet[Theorem~1.8 and Corollary~1.9]{FF_CharlesPoonen2016}.
For a closed point of degree greater than
\(e\deg\operatorname{div}_\infty h+
d\deg\operatorname{div}_\infty g\), the
resultant
\[
        \operatorname{Res}_Z
          (Z^d-\zeta_xh,\ (Z-1)^e-\eta_xg)
\]
vanishes identically, by its uniform pole
bound. Its coefficients lie in \(K\),
since all roots of unity do. Restriction
injectivity makes it vanish in \(K(X)\).
Its common root \(z\) is rational: a
conjugate satisfies \(z'=\alpha z\) and
\(z'-1=\beta(z-1)\), with \(\alpha,\beta\)
constant roots of unity, and nonconstancy
forces \(\alpha=\beta=1\). It is a unit
because \(z^d=\zeta_xh\).
Equation (FF.35.6) then gives
\(\lambda_0=\kappa_p(z)\) geometrically
with rational coefficients. Rational
coefficients suffice for this division
by \(d\), and a rationally zero unit
Kummer class still has zero divisor.
Normalize \(z(x_0)=b_0\). Evaluation at
\(x_0\) kills the constant cohomology
difference, so (FF.35.3) becomes
\(\kappa_p(b(x))=\kappa_p(z(x))\)
at every finite-extension point. Hence
\(b(x)/z(x)\in\mu_\infty\subset K\).

If \(\lambda_0=0\), (FF.35.3) puts all
values \(b(x)/b_0\) in \(\mu_\infty\subset K\).
Choose \(\beta\) in a proper finite
extension of \(K\), and apply the same
construction to \(T-\beta\). Such an
extension exists by Step~4. If its
weight-zero class is nonzero, the preceding
resultant gives a nonconstant unit, whereas
all values of that unit lie in the fixed
finite base extension containing \(\beta\):
indeed \(b(x)\in K\), and its root-of-unity
ratio also belongs to \(K\). This holds for
all geometric points, including points defined
over arbitrarily larger finite fields: every
\(b(x)-\beta\) still lies in \(K(\beta)\).
The
large-degree-point argument in Step~5
of Claim~\ref{clm:FFstrictreconstruct}
excludes this. If its weight-zero class
is zero, comparing the two root-of-unity
ratios and using the linear independence
of \(1,\beta\) over \(K\) makes
\(b(x)=b_0\). Thus the zero case has a
constant actual point function.

Finally, for a nonconstant weight-zero
class, (FF.35.3) and its realizing unit
give the same constant-scalar ambiguity
as in Step~6 of
Claim~\ref{clm:FFstrictreconstruct}.
Apply the construction also to \(T-1\),
and use a closed point of degree larger
than the relevant pole bound and the
degree of the nonconstant unit as a map
to \(\mathbb P^1\). Linear independence
at that point forces the two units and
the actual values to have one fixed
scalar relation. Applying \(T-a\) deals
with the remaining base-field-value
points. Each nonconstant point function
remains nonconstant on its inverse-image
open: otherwise the inverse image of a
closed target point would contain a dense
open of the integral source and hence
be the whole source. In higher dimension,
the same argument is made on a smooth
geometrically integral curve through the
base point and the point being tested,
using the cited Charles--Poonen theorem.
These are precisely the algebraic
scalar-removal steps already proved there;
they require only (FF.35.3), all roots of
unity in the base, and Step~4.

The weight decomposition used throughout
is compatible with restriction to opens,
finite scalar extensions, and pullback
along the curves: it is the eigenspace
decomposition of the normalized inertia
Lie generator. In particular a geometric
Kummer functional of an algebraic unit
belongs to weight zero, and comparisons
of such units in this last step preserve
the normalized class \(\lambda_0\).
We have therefore obtained a unit whose
actual geometric point function is exactly
that of \(\varphi\). Finite descent returns
it to the original \(K\). The point-function
criterion cited at the start proves
\(\mathcal R(K)\).
Every subfield \(F\subset K\) has
\(K/F\) algebraic, since \(K/\mathbb Q\)
is algebraic. Algebraic descent, Claim~\ref{clm:FFdescent},
proves \(\mathcal R(F)\).
\end{proof}
